\begin{titlepage}
	\centering
	
	% Definimos los colores locales para el TikZ
	\colorlet{colorRbase}{blue!10}       
	\colorlet{colorRline}{blue!50}       
	\colorlet{colorRcell}{blue!40}       
	\colorlet{colorRtext}{black}         
	\colorlet{colorRhatBase}{gray!15}    
	\colorlet{colorRhatLine}{gray!50}    
	\colorlet{colorRhatPred}{NaranjaTexto}    
	\colorlet{colorProj}{VerdeDecorado}        
	
	\vspace*{\fill} % Empuja el contenido hacia abajo para centrar en la hoja
	
	\begin{tikzpicture}
		\node[font=\fuentedeportada\fontsize{60}{45}\selectfont, text=black, align=center] at (0,0){Sistemas de \\ Recomendación};
	\end{tikzpicture}
	
	\vfill % IMPORTANTE: Las líneas en blanco antes y después son necesarias
	
	\begin{tikzpicture}[
		% Configuramos la perspectiva isométrica personalizada
		x={(-0.866cm,-0.5cm)}, 
		y={(0.866cm,-0.5cm)}, 
		z={(0cm,1cm)},
		scale=1,
		font=\fuentedeportada
		]
		
		% VARIABLES PARA LUZ EN ÁNGULO RASANTE
		\def\dx{-4}   
		\def\dy{2} 
		
		% ==========================================
		% 1. CAPA INFERIOR: MATRIZ COMPLETADA
		% ==========================================
		\fill[colorRhatBase] (\dx,\dy,0) -- (5+\dx,\dy,0) -- (5+\dx,5+\dy,0) -- (\dx,5+\dy,0) -- cycle;
		
		%Grilla
		\foreach \i in {0,...,5} {
			\draw[black,thick] (\i+\dx,\dy,0) -- (\i+\dx,5+\dy,0);
			\draw[black,thick] (\dx,\i+\dy,0) -- (5+\dx,\i+\dy,0);
		}
		
		\foreach \x/\y/\v in {
			0/0/5, 1/0/2, 2/0/3, 3/0/4, 4/0/2,
			0/1/4, 1/1/3, 2/1/4, 3/1/1, 4/1/5,
			0/2/1, 1/2/5, 2/2/2, 3/2/4, 4/2/3,
			0/3/3, 1/3/4, 2/3/5, 3/3/3, 4/3/1,
			0/4/2, 1/4/1, 2/4/3, 3/4/4, 4/4/5%
		} {
			\node[text=colorRhatPred] at (\x+0.5+\dx, \y+0.5+\dy, 0) {\v};
		}
		
		\foreach \x/\y/\v in {
			0/0/5, 2/1/4, 1/2/5, 4/3/1, 2/4/3
		} {
			\node[text=black, font=\fuentedeportada] at (\x+0.5+\dx, \y+0.5+\dy, 0) {\v};
		}
		
		% ==========================================
		% 2. RAYOS DE PROYECCIÓN
		% ==========================================
		\foreach \x/\y in {0/0, 5/0, 0/5, 5/5} {
			\draw[dashed, thick, colorProj] (\x,\y,6) -- (\x+\dx,\y+\dy,0);
		}
		
		% ==========================================
		% 3. CAPA SUPERIOR: MATRIZ INCOMPLETA
		% ==========================================
		\fill[colorRbase, opacity=0.6] (0,0,6) -- (5,0,6) -- (5,5,6) -- (0,5,6) -- cycle;
		
		\foreach \i in {0,...,5} {
			\draw[black, thick] (\i,0,6) -- (\i,5,6);
			\draw[black, thick] (0,\i,6) -- (5,\i,6);
		}
		
		\foreach \x/\y/\v in {
			0/0/5, 2/1/4, 1/2/5, 4/3/1, 2/4/3
		} {
			\fill[colorRcell, opacity=0.8] (\x,\y,6) -- (\x+1,\y,6) -- (\x+1,\y+1,6) -- (\x,\y+1,6) -- cycle;
			\draw[colorRline, thick] (\x,\y,6) -- (\x+1,\y,6) -- (\x+1,\y+1,6) -- (\x,\y+1,6) -- cycle;
			\node[text=colorRtext, font=\fuentedeportada] at (\x+0.5, \y+0.5, 6) {\v};
		}
		
		% ==========================================
		% 4. ETIQUETAS
		% ==========================================
		\foreach \c in {1,...,5} {
			\node[anchor=south east, text=black] at (\c-0.5, 0, 6) {$i_{\c}$};
			\node[anchor=south west, text=black] at (0, \c-0.5, 6) {$u_{\c}$};
		}
	\end{tikzpicture}
	
	\vfill % Separación vertical para empujar los nombres hacia abajo
	
	\begin{tikzpicture}
		\node[font=\fuentedeportada\fontsize{20}{25}\selectfont, text=black, align=center] at (0,0){Arias, Blanc, Comelli};
	\end{tikzpicture}
	
	\vspace*{\fill} % Empuja el contenido hacia arriba para equilibrar la hoja
\end{titlepage}